\section{Modelling Methodology}
\label{sec:method}

\subsection{Model family budget and pipelines}
\label{sec:families}

The brief allows two to three families and asks for depth rather than breadth.
We use three, reused across both task types, each arranged as a
\emph{complexity ladder} whose rung 0 is the family's simplest variant --- the
only way ``the added complexity was worth it'' becomes a finding rather than an
assumption. \textbf{A --- Linear} runs \code{LinearRegression} /
\code{LogisticRegression} at rung 0 and \code{Ridge} / $\ell_2$-penalised
logistic at rung 1. \textbf{B --- Tree-based} runs \code{DecisionTree} $\to$
\code{RandomForest} $\to$ \code{HistGradientBoosting}, the three steps
Section~\ref{sec:lit} argued from the literature. \textbf{C --- Ensemble} is
soft \code{Voting} and \code{Stacking} assembled from the winners of A and B;
its members are complete pipelines rather than bare estimators, because the two
families need different encodings and one shared preprocessor would force a
single choice on both. The full specification, the search spaces and the
selected values of every tuned model are in Appendix~\ref{app:hyperparameters}.

Every model is a \code{scikit-learn} \code{Pipeline}: median imputation, one-hot
encoding of the four categoricals with rare levels folded into one bucket, then
--- for the linear family only --- standardisation, since trees are invariant to
monotone rescaling and a shared penalty is not \citep{pedregosa2011sklearn}. The
linear family also drops one level per categorical: Phase 1 called least squares
``rank-deficient, so its coefficients explode'', and re-running it four ways
(Appendix~\ref{app:encoding}) shows the predictions were always fine and only the
coefficients were not; with full rank restored, Ridge and OLS agree to three
decimals, so regularisation buys numerical stability here, not accuracy.

\subsection{Cross-validation, and how a model gets chosen}
\label{sec:cv}

Delivery performance drifts hard --- mean lead time falls \textbf{5.34 days}
between training and test --- so shuffled $k$-fold lets a model learn from the
future \citep{bergmeir2012cv}. Hyperparameters are tuned by \code{TimeSeriesSplit}
with five expanding windows inside the training period and
\code{RandomizedSearchCV} \citep{bergstra2012random} over 30 draws per model (six
for the costly random forests), scoring MAE or average precision.

\emph{Selecting} between tuned models is a separate question, and both obvious
answers fail here. Averaging cross-validated scores fails because the folds differ
enormously in difficulty --- lead time runs 11.5--16.5 days across the folds
against 8.3 in the test window --- so the \emph{ordering} does not survive:
cross-validated MAE ranks Ridge above the boosted trees, while validation and test
rank them the other way round (Figure~\ref{fig:ladder}b). The validation period
alone fails too: its two months pull in opposite directions, 8.2\% late in May
2018 and \textbf{1.4\%} in June. Each candidate is therefore scored on all six
out-of-sample windows, \textbf{ranked within each}, and the winner is the best
mean rank over the \textbf{three most recent} --- ranks because only ranks are
comparable across windows of different difficulty, three because they cover
roughly February to June 2018 and neither odd month decides alone. All six windows
are reported for every candidate in Appendix~\ref{app:protocol}; a shuffled,
customer-disjoint \code{StratifiedGroupKFold} is computed only as a contrast
(Figure~\ref{fig:protocol}). The
selected models are \emph{not} refitted on train plus validation before the test
window is scored, so what is reported is exactly what is deployed.

\subsection{Imbalance, metrics aligned to the action, and reproducibility}
\label{sec:imbalance}

Both targets are imbalanced --- 7.4\% late and 11.0\% low review in the test
period --- so accuracy alone is meaningless and every results row carries
\code{majority\_accuracy} beside \code{accuracy}, with PR-AUC reported against
its no-skill floor \citep{saito2015prplot}. Following
\citet{goorbergh2022imbalance} we fit on the observed prior and move the
\emph{decision threshold} instead, tuned on the validation period by maximising
F1, and Section~\ref{sec:results-review} tests the alternative rather than
assuming the result; no resampling is used anywhere \citep{chawla2002smote}.
Calibrated probabilities are a requirement rather than a nicety, because
Section~\ref{sec:cost} multiplies them by a cost \citep{elkan2001cost}; they are
reported, not assumed \citep{niculescu2005calibration}.

A second alignment follows from Table~\ref{tab:levers}. PR-AUC summarises the
whole ranking and is the right criterion for \emph{choosing} between models,
but neither team acts on the whole ranking: customer experience contacts the
$K$ riskiest customers its agents can reach in a day, and a warning, where one
is shown, goes to the riskiest orders. Section~\ref{sec:results-capacity} therefore reports
Precision@$K$ and Recall@$K$ for both flags \citep{provost2013data,
verbeke2012churn}, with $K$ as a share of the day's volume so that the same
table holds on a quiet Tuesday and on Black Friday, and Section~\ref{sec:cost}
chooses the operating point as a capacity rather than as a threshold.

\code{SEED = 42} throughout. \code{src/train\_models.py} fits everything and
writes the models, metadata and result tables, \code{src/customer\_view.py} the
customer counts; \code{src/phase3\_stats.py}
distils those into \code{report/stats\_phase3.txt}, and no number enters this
report that it cannot regenerate; every figure is drawn from the same tables.
Artifacts record the \code{scikit-learn} version that wrote them and the loader
checks it; \code{tests/test\_leakage.py} runs \textbf{\checkcount{}} validation
checks, including one that ties the deployed artifacts to the metadata and to
this text (Appendix~\ref{app:repro}).
